%!TEX root = ../main.tex

\chapter{Knowledge Graph Construction}\label{ch:kg}

This chapter describes the construction of a \gls{kg} from the \gls{mdd} corpus,
the first stage beyond passage retrieval. The work is in progress: the graph
design, the extraction pipeline and its quality checks are complete on a set of
pilot documents, while the extraction over the full corpus is deferred until the
manual precision annotation described in Section~\ref{sec:kg:status} is
available. Section~\ref{sec:kg:design} presents the three-layer design.
Sections~\ref{sec:kg:structure}--\ref{sec:kg:extraction} cover the structural
layer, the schema and the extraction pipeline. Sections~\ref{sec:kg:quality}
and~\ref{sec:kg:er} report the quality checks and entity resolution, and
Section~\ref{sec:kg:status} summarises the state of the work.


% ===========================================================================
\section{Design}\label{sec:kg:design}

The construction follows the methodology of \textcite{negro2026kg} (chapters 5
and 6), which separates a \gls{kg} into three layers:

\begin{enumerate}
    \item a \emph{structural layer} that mirrors the documents themselves (domains,
          documents, sections and spans), built without an \gls{llm};
    \item a \emph{metagraph} of raw mentions: every entity and relation an
          \gls{llm} extracts from a section, with a verbatim evidence quote and a link
          to that section;
    \item a \emph{\gls{kg} layer} of normalised entities, obtained by resolving
          mentions that refer to the same real-world entity.
\end{enumerate}

Keeping raw mentions separate from resolved entities makes every fact traceable
to its source text and allows entity resolution to be revised without
re-extracting. The graph is stored in Neo4j~\cite{neo4j} (version~5, Community
Edition).


% ===========================================================================
\section{Structural Layer}\label{sec:kg:structure}

The structural layer is loaded deterministically from the corpus:
\begin{center}
\small\ttfamily
(:Domain)-[:HAS\_DOCUMENT]->(:Document)-[:HAS\_SECTION]->(:Section)-[:HAS\_SPAN]->(:Span)\\
(:Section)-[:NEXT]->(:Section)
\end{center}
Sections are the retrieval units of Chapter~\ref{ch:rag} and use the same
identifiers, so a retrieved section maps directly to its graph node; span
identifiers likewise match the dataset's grounding annotations. All writes use
\texttt{MERGE} on these identifiers, so the load is idempotent.
Table~\ref{tab:kg:structure} gives the resulting counts.

\begin{table}[ht]
    \centering
    \caption{Structural layer of the \gls{kg}.}
    \label{tab:kg:structure}
    \begin{tabular}{lr@{\hspace{3em}}lr}
        \toprule
        Node label & Count & Relationship & Count \\
        \midrule
        Domain   & 4        & HAS\_DOCUMENT & 488 \\
        Document & 488      & HAS\_SECTION  & 4{,}114 \\
        Section  & 4{,}114  & HAS\_SPAN     & 35{,}659 \\
        Span     & 35{,}659 & NEXT          & 3{,}626 \\
        \bottomrule
    \end{tabular}
\end{table}

Two section flags support extraction planning. The corpus HTML contains no
\texttt{<table>} elements (they were removed when the dataset was built), so
\texttt{is\_table} is a text heuristic: a section is flagged when one of its spans
enumerates at least six short comma- or semicolon-separated items, such as state
codes or vehicle registration classes (37 sections). \texttt{is\_boilerplate}
flags sections whose text occurs in ten or more documents, such as navigation
widgets and ``Related PDFs:'' headers (300 sections).


% ===========================================================================
\section{Schema Design}\label{sec:kg:schema}

\subsection{Open Extraction Pilot}

The schema was derived from data rather than fixed in advance. An open-extraction
pilot asked the \gls{llm} to extract entities and relations with freely chosen
types from 32 sections (eight per domain, sampled with a fixed seed), each with a
verbatim evidence quote. It produced 318 entities and 264 relations, but with 139
distinct entity types and 180 distinct relation types, roughly one new type for
every two mentions. Open extraction alone therefore yields no usable schema.

The pilot also exposed systematic errors that later became regression tests:
relations in the wrong direction (e.g.\ ``supervised driving practice
\textsc{requires} road test''), relation names that encode their target
(\texttt{REQUIRES\_VISION\_TEST}), self-loops, and entity names written as CamelCase
identifiers instead of the text's wording (69\% of names). Evidence quotes were
reliable: 97.4\% matched the section text exactly or up to punctuation, and
manual inspection of the remaining 15 found no invented content.

\subsection{Closed Schema}

Every open type was mapped by hand onto a candidate schema, merged, or dropped.
Dropped types were mostly page interface elements, statistics, table rows and
calculation details. The resulting closed schema has nine entity types
(Table~\ref{tab:kg:entities}) and eleven relation types with fixed directions
(Table~\ref{tab:kg:relations}); it covers 78.9\% of the pilot's entity mentions
and 76.1\% of its relations. Amount, TimeLimit and ContactPoint carry a verbatim
\texttt{value} (e.g.\ ``\$25.00'', ``every 12 months'', a phone number), because
the pilot showed that such values otherwise appear only in the evidence and cannot
answer ``how much'' or ``how long'' questions. Requirement and Outcome keep the
full condition or consequence as text.

\begin{table}[ht]
    \centering
    \small
    \caption{Entity types of the closed schema.}
    \label{tab:kg:entities}
    \begin{tabular}{l>{\raggedright\arraybackslash}p{0.62\textwidth}}
        \toprule
        Type & Definition \\
        \midrule
        Service      & something a citizen can apply for, receive or complete (benefit, programme, licence, loan, test, appeal, \ldots) \\
        Organization & an agency, office, school, servicer or other body that provides, issues or decides \\
        Document     & a form, certificate, proof, card, letter or notice \\
        Requirement  & a condition or obligation: eligibility criterion, item to present, restriction \\
        Amount       & a fee, tax, premium, benefit amount or limit, with its value \\
        TimeLimit    & a deadline, duration, frequency or waiting period, with its value \\
        Role         & a category of person the content applies to \\
        Outcome      & a result or consequence of an action or condition \\
        ContactPoint & a phone number, website, account, office or information page \\
        \bottomrule
    \end{tabular}
\end{table}

\begin{table}[ht]
    \centering
    \small
    \caption{Relation types of the closed schema, with allowed source and target
             types. ALTERNATIVE\_TO is symmetric and stored in one canonical
             direction.}
    \label{tab:kg:relations}
    \begin{tabular}{l>{\raggedright\arraybackslash}p{0.36\textwidth}>{\raggedright\arraybackslash}p{0.33\textwidth}}
        \toprule
        Relation & Source & Target \\
        \midrule
        REQUIRES       & Service, Document, Role; Requirement (to Document only) & Requirement, Document, Service \\
        PROVIDED\_BY   & Service, Document, Amount & Organization \\
        HAS\_AMOUNT    & Service, Document & Amount \\
        HAS\_TIME\_LIMIT & Service, Requirement, Document & TimeLimit \\
        APPLIES\_TO    & Service, Requirement, Amount, TimeLimit & Role, Service \\
        LEADS\_TO      & Service, Requirement, Outcome & Outcome, Service \\
        EXEMPT\_FROM   & Role, Service, Document & Requirement, Amount, Service \\
        DETERMINED\_BY & Amount, Service, Outcome & Requirement, Role, Amount \\
        AVAILABLE\_AT  & Service, Document; Organization (to ContactPoint only) & ContactPoint, Organization \\
        PART\_OF       & Service, Document, Amount & Service, Document, Amount \\
        ALTERNATIVE\_TO & Service, Document, Requirement & Service, Document, Requirement \\
        \bottomrule
    \end{tabular}
\end{table}


% ===========================================================================
\section{Extraction Pipeline}\label{sec:kg:extraction}

Each extraction call gives the \gls{llm} the domain, document title, section
heading, section text and the schema with its source and target rules, and
requests a strict JSON-schema output. The prompt requires names in the text's own
wording, an optional abbreviation as \texttt{alias} (e.g.\ ``Department of Motor
Vehicles'' / ``DMV''), one contiguous verbatim evidence quote per item, and at most
15 entities per call. Post-processing normalises names, removes duplicate entities,
self-loops and relations whose endpoints were not extracted, stores symmetric
relations in one direction, and flags relations that violate the source and target
rules; flagged relations are kept in the raw output but not loaded into the graph.

Extraction is planned in \emph{units} (Table~\ref{tab:kg:plan}). Sections shorter
than 40 characters and boilerplate sections are skipped; a text repeated in several
sections is extracted once and attached to every section where it occurs, so each
section keeps its own mention links; and texts longer than 1{,}500 characters are
split into chunks at sentence boundaries, each with its own entity cap, while
evidence is still checked against the whole section.

\begin{table}[ht]
    \centering
    \caption{Extraction plan for the full corpus.}
    \label{tab:kg:plan}
    \begin{tabular}{lr}
        \toprule
        & Sections / calls \\
        \midrule
        Sections in the corpus             & 4{,}114 \\
        \quad skipped as boilerplate       & 300 \\
        \quad skipped as shorter than 40 characters & 262 \\
        \quad extracted                    & 3{,}552 \\
        Distinct texts (extraction units)  & 3{,}194 \\
        Units split into chunks            & 238 \\
        \midrule
        \gls{llm} calls                    & 3{,}542 \\
        \bottomrule
    \end{tabular}
\end{table}


% ===========================================================================
\section{Quality Checks}\label{sec:kg:quality}

The closed-schema pipeline was evaluated on every section of eight whole documents
(two per domain, 90 sections) before any full-corpus run.
Table~\ref{tab:kg:pilots} compares the open pilot with two closed-schema runs.

\begin{table}[ht]
    \centering
    \small
    \caption{Extraction quality across pilots. The open pilot used 32 sections; the
             closed-schema runs used the 90 sections of eight whole documents
             (pilot 4 with pinned sampling parameters and the ContactPoint value
             rule).}
    \label{tab:kg:pilots}
    \begin{tabular}{lrrr}
        \toprule
        & Open pilot & Closed, pilot 2 & Closed, pilot 4 \\
        \midrule
        Entities / relations              & 318 / 264 & 511 / 368 & 509 / 372 \\
        Distinct entity / relation types  & 139 / 180 & 9 / 11 & 9 / 11 \\
        CamelCase entity names            & 69\%  & 1.4\% & 1.2\% \\
        Evidence quoted verbatim          & 81.1\% & 96.8\% & 96.0\% \\
        Relations violating type rules    & ---   & 8.4\% & 15.1\% \\
        Self-loops (removed)              & 7     & 0 & 7 \\
        \bottomrule
    \end{tabular}
\end{table}

Most evidence quotes that fail the verbatim check are caused by the corpus itself,
which contains run-together words (``areunder'', ``yourbenefits'') that the
\gls{llm} silently corrects. Relations that violate the type rules are not
dominated by any single pattern and are dropped at load time. The known errors of
the open pilot were turned into 17 regression tests that run against cached model
output; the reversed \textsc{requires} relations and object-specific relation names
are fixed, while four documented failures remain (an exemption still extracted in
the reverse direction, a fee-table enumeration, 1.4\% raw self-loops, and one
stitched quote).

\subsection{Run-to-Run Stability}

Running the same pilot twice revealed that the extraction is far from
deterministic: only 48\% of the triples of one run reappear in the other (triple
Jaccard 0.33), and about a quarter of the entities differ, mainly in granularity
(one entity or two) and in which peripheral facts are extracted. The deployment
offers no deterministic setting: the \texttt{seed} parameter is not supported, and
temperature is accepted only when reasoning effort is set to ``none''.
Table~\ref{tab:kg:determinism} compares the available configurations.

\begin{table}[ht]
    \centering
    \caption{Determinism of the extraction under the available sampling
             configurations (four sections, five repetitions each; mean pairwise
             Jaccard of the extracted triples).}
    \label{tab:kg:determinism}
    \begin{tabular}{lr}
        \toprule
        Configuration & Triple Jaccard \\
        \midrule
        temperature 0, reasoning effort ``none'' & 0.34 \\
        temperature 0, top-$p$ 0                 & 0.27 \\
        reasoning effort ``medium''              & 0.24 \\
        reasoning effort ``low''                 & 0.17 \\
        top-$p$ 0                                & 0.12 \\
        \bottomrule
    \end{tabular}
\end{table}

The most stable configuration is the default one, which is now pinned explicitly;
a repeated pilot with it reproduced 44\% of the triples (Jaccard 0.28), within the
noise of the earlier measurement. A single extraction pass therefore misses facts
that a second pass finds, and aggregating several runs (union for recall,
intersection for precision) is the remaining lever.


% ===========================================================================
\section{Entity Resolution}\label{sec:kg:er}

Entity resolution merges mentions into the entities of the \gls{kg} layer. Mentions
are first blocked by type: Service, Requirement, Role and Outcome are compared only
within a domain, while Organization, Document, ContactPoint, Amount and TimeLimit
may be merged across domains. Mentions extracted with conflicting types are assigned
the majority type. Candidate pairs come from three sources: deterministic rules
(abbreviation aliases, acronyms, and shared form numbers such as IRP-6), the cosine
similarity of name embeddings (\texttt{bge-small-en-v1.5}~\cite{bge2023}), and graph
signals that raise confidence when two mentions share a document or relation
neighbours, as in~\textcite[ch.~6.1.4]{negro2026kg}.

\subsection{Threshold Calibration}

The \gls{llm} was asked to judge every candidate pair scoring at least 0.75 and a
random sample below, to calibrate the similarity thresholds
(Table~\ref{tab:kg:calibration}).

\begin{table}[ht]
    \centering
    \caption{Share of candidate pairs judged to be the same entity, by similarity
             score (eight pilot documents).}
    \label{tab:kg:calibration}
    \begin{tabular}{lrr}
        \toprule
        Score & Pairs & Judged same \\
        \midrule
        0.65--0.75 (sample) & 100 & 1\% \\
        0.75--0.80 & 199 & 9\% \\
        0.80--0.85 & 117 & 10\% \\
        0.85--0.90 &  72 & 25\% \\
        0.90--0.95 &  48 & 44\% \\
        $\geq 0.95$ &  19 & 53\% \\
        \bottomrule
    \end{tabular}
\end{table}

No similarity score is safe for automatic merging: even above 0.95 about half of the
pairs are distinct entities whose names differ in a single qualifier, such as
subsidised and unsubsidised loans or medical and non-medical hearings. Only rule
matches are therefore merged automatically. Pairs below 0.80 are never merged: a
manual reading of the 0.75--0.80 band found that almost all positive judgments there
were wrong merges. Every pair in between is confirmed by the \gls{llm}.

Three safeguards were added after a first run produced chained clusters (A merged
with B and B with C, although A and C were judged different): two clusters are not
merged if any pair across them was rejected; Amount, TimeLimit and ContactPoint
groups with different stated values are never merged; and an alias counts as merge
evidence only if it looks like an abbreviation and is unique in its block, because
the \gls{llm} occasionally fills the alias field with a category label.

\subsection{Batched Confirmation and Cost}

Confirming one pair per call does not scale: at corpus size, blocks become dense
enough that almost every name group has at least five candidates, so even with a
cap of five candidates per group, pairwise confirmation would need about 66{,}000
calls. Batched confirmation sends one query entity with up to five candidates per
call and asks for a separate decision per candidate, preceded by an explicit
statement of what distinguishes them. On the pilot it reduced the calls from 371 to
162 at the same acceptance rate (17\% against 18\%), and its decisions agree with
pairwise confirmation on 85\% of pairs; since the pairwise judge agrees with itself
on 98\% of re-asked pairs, the remaining difference is a genuine change in behaviour,
which the manual cluster review will arbitrate. A simulation calibrated on the pilot
estimates about 21{,}000 batched calls for the full corpus. With the final settings
(batched confirmation, lower threshold 0.80) the pilot's 508 mentions resolve into
382 entities.


% ===========================================================================
\section{Status}\label{sec:kg:status}

The structural layer is complete in Neo4j. The extraction pipeline, its schema and
the entity resolution procedure are implemented and calibrated on eight pilot
documents. Before the full extraction of 3{,}542 calls, a manual annotation of 140
extracted triples from three of the pilot documents (one each from DMV, SSA and VA),
together with the triples the extraction missed, will measure precision and recall
per relation and entity type. Precision is also reported separately for the triples
that both independent runs extracted, which indicates whether aggregating runs would
improve precision. The proposed entity merges are reviewed in the same way. These
measurements decide whether to run the full extraction as designed, with several
aggregated runs, or with a revised prompt.
